\begin{DefinitionBox}{Geradengleichung}{real_functions_affine}
  Eine \emph{affine Funktion} \EnglishTerm{affine function} besitzt mit
  $m,b\in\R$ und $D=\R$ die Gleichung
  \[
    f(x)=mx+b.
  \]
  Ihr Graph ist eine Gerade. Dabei ist $m$ die \emph{Steigung}
  \EnglishTerm{slope}; $b=f(0)$ ist der $y$-Achsenabschnitt.
  Für $b=0$ ist die Zuordnung proportional.
\end{DefinitionBox}
In der Schulmathematik nennt man $f(x)=mx+b$ oft eine \emph{lineare Funktion}
\EnglishTerm{linear function}. Im strengeren algebraischen Sinn ist nur
$f(x)=mx$ linear. Eine senkrechte Gerade $x=c$ hat keine endliche Steigung
und ist kein Graph einer Funktion $y=f(x)$.


\begin{KeyPointBox}{Steigung und Lage der Geraden}{real_functions_line_parameters}
  Das Vorzeichen von $m$ bestimmt die Richtung des Graphen:
  \[
    \begin{array}{c@{\qquad}c@{\qquad}c}
      m>0&m=0&m<0\\
      \text{steigende Gerade}&\text{waagerechte Gerade}&\text{fallende Gerade}
    \end{array}
  \]
\end{KeyPointBox}

\clearpage
\subsection{Nullstelle einer Geraden}
\label{subsec:real_functions_line_zero}
\begin{KeyPointBox}{Nullstelle und $y$-Achsenabschnitt einer Geraden}{real_functions_line_intercepts}
  Für eine Gerade $f(x)=mx+b$ mit $m,b\in\R$ ergeben sich die beiden
  Achsenabschnitte direkt aus der Geradengleichung:
  \begin{itemize}[leftmargin=1.3em,itemsep=0.4em,topsep=0.3em]
    \item $y$-Achsenabschnitt:\par
      \[y_0=f(0)=b.\]
    \item Nullstelle für $m\ne0$:\par
      \[f(x_0)=0\Longleftrightarrow mx_0+b=0
        \Longleftrightarrow x_0=-\frac{b}{m}.\]
  \end{itemize}
  Die zugehörigen Punkte sind $(0,b)$ und $(x_0,0)$.
  Für $m=0$ gilt: Ist $b\ne0$, gibt es keine Nullstelle; ist $b=0$,
  ist jede reelle Zahl eine Nullstelle.
\end{KeyPointBox}
\begin{ExampleBox}{Nullstelle und $y$-Achsenabschnitt einer Geraden}{real_functions_line_intercepts_example}
  Wir betrachten die Geradengleichung $f(x)=mx+b$ mit $m=\frac12$ und $b=1$.
  \begin{minipage}[c]{0.50\linewidth}
     Mit den Formeln aus
    Merke~\ref{key:real_functions_line_intercepts} berechnen wir:
    \begin{itemize}[leftmargin=1.3em,itemsep=0.4em,topsep=0.3em]
      \item $y$-Achsenabschnitt:\par
        $y_0=b=\Cloze{1.2cm}{1}$.
      \item Nullstelle:\par
        $x_0=-\frac{b}{m}
            =-\frac{1}{\frac12}
            =\Cloze{1.2cm}{-2}$.
    \end{itemize}
    Die Achsenpunkte sind $(0,\Cloze{1.2cm}{1})$ und
    $(\Cloze{1.2cm}{-2},0)$.
  \end{minipage}\hfill
  \begin{minipage}[c]{0.49\linewidth}
    \raggedleft
    \CachedTikzFigure{content/resources/chapters/real_functions/tikz/functions_line_intercepts.tex}
  \end{minipage}
\end{ExampleBox}

\clearpage
\subsection{Eine Gerade aus zwei Punkten bestimmen}
\label{subsec:real_functions_two_points}
Zwei Punkte $P_A=(x_A,y_A)$ und $P_B=(x_B,y_B)$ mit $x_A\ne x_B$
bestimmen genau eine Gerade. Die Gleichung $y=mx+b$ enthält die beiden
Unbekannten $m$ und $b$; jeder Punkt liefert dafür eine Gleichung.
So entsteht ein quadratisches LGS.
Vor der erweiterten Matrix ordnen wir die Terme nach $b$ und $m$:
\[
  \begin{aligned}
    mx_A+b&=y_A,\\
    mx_B+b&=y_B
  \end{aligned}
  \quad\Longleftrightarrow\quad
  \begin{aligned}
    b+x_A m&=y_A,\\
    b+x_B m&=y_B
  \end{aligned}
  \quad\Longleftrightarrow\quad
  \left[\begin{array}{@{}cc|c@{}}
    \Cloze{0.7cm}{1}&\Cloze{0.8cm}{x_A}&\Cloze{0.8cm}{y_A}\\[0.3em]
    \Cloze{0.7cm}{1}&\Cloze{0.8cm}{x_B}&\Cloze{0.8cm}{y_B}
  \end{array}\right].
\]
Wie in Kapitel 2 eliminieren wir zunächst einen Koeffizienten und setzen
anschließend von der letzten Zeile aus rückwärts ein:
\par\noindent\textbf{Vorwärtselimination:} Die Operation steht rechts neben der Matrix,
auf die sie angewendet wird:
\begin{ClozeArea}[34]{6}
  \[
    \begin{array}{r@{\;}r@{\qquad}l}
      &\left[\begin{array}{@{}cc|c@{}}
        1&x_A&y_A\\
        1&x_B&y_B
      \end{array}\right]
      &\raisebox{-0.65em}{\(\TransformOp{Z_2-Z_1}\)}\\[1.2em]
      \Longleftrightarrow
      &\left[\begin{array}{@{}cc|c@{}}
        1&x_A&y_A\\
        \textcolor{THMBlue}{0}&x_B-x_A&y_B-y_A
      \end{array}\right]
      &\text{Zeilenstufenform}
    \end{array}
  \]
\end{ClozeArea}
\noindent\textbf{Rückwärtseinsetzen:} Aus Zeile 2 folgt die Steigung:
\begin{ClozeArea}[34]{4}
  \[
    (x_B-x_A)m=y_B-y_A
    \quad\Longleftrightarrow\quad
    m=\frac{y_B-y_A}{x_B-x_A}.
  \]
\end{ClozeArea}
\noindent Einsetzen in Zeile 1 liefert den $y$-Achsenabschnitt:
\begin{ClozeArea}[34]{5}
  \[
    b=y_A-x_A m
     =y_A-x_A\frac{y_B-y_A}{x_B-x_A}
     =\frac{x_By_A-x_Ay_B}{x_B-x_A}.
  \]
\end{ClozeArea}
\begin{KeyPointBox}{Zweipunkteform}{real_functions_two_point_form}
  \begin{minipage}[c]{0.52\linewidth}
    Die Gerade durch $P_A=(x_A,y_A)$ und $P_B=(x_B,y_B)$ mit
    $x_A\ne x_B$ hat die Gleichung
    \[
      f(x)=mx+b
    \]
    mit
    \[
      \begin{aligned}
        m&=\frac{y_B-y_A}{x_B-x_A}
          \quad\textcolor{green!55!black}{\text{(Steigungsdreieck)}}\\[0.4em]
        b&=y_A-mx_A
          =\frac{x_By_A-x_Ay_B}{x_B-x_A}.
      \end{aligned}
    \]
  \end{minipage}\hfill
  \begin{minipage}[c]{0.46\linewidth}
    \centering
    \CachedTikzFigure{content/resources/chapters/real_functions/tikz/functions_two_point_line.tex}
  \end{minipage}
\end{KeyPointBox}
\begin{ExampleBox}{Gerade durch zwei Punkte}{real_functions_two_point_numeric_example}
  Welche Gerade verläuft durch $P_A=(1,1)$ und $P_B=(3,2)$?
  \par\smallskip
  \begin{minipage}[c]{0.45\linewidth}
    Aus den beiden Punkten bestimmen wir zunächst die Steigung und danach
    den $y$-Achsenabschnitt:
    \[
      \begin{aligned}
        m&=\frac{2-1}{3-1}=\Cloze{1.4cm}{\frac12},\\[0.5em]
        b&=1-m\cdot1=\Cloze{1.4cm}{\frac12}.
      \end{aligned}
    \]
    Damit lautet die Geradengleichung
    \[
      f(x)=\Cloze{3.5cm}{\frac12x+\frac12}.
    \]
  \end{minipage}\hfill
  \begin{minipage}[c]{0.54\linewidth}
    \raggedleft
    \CachedTikzFigure{content/resources/chapters/real_functions/tikz/functions_two_point_numeric_example.tex}
  \end{minipage}
\end{ExampleBox}
\begin{ExampleBox}{Kalibrierung eines Wegsensors}{real_functions_sensor_calibration}
  Ein Sensor liefert bei $s=10\,\mathrm{mm}$ die Spannung
  $U=1\,\mathrm{V}$ und bei $s=50\,\mathrm{mm}$ die Spannung
  $U=9\,\mathrm{V}$. Wir nehmen einen affinen Zusammenhang an und
  bestimmen $U(s)$ sowie die Spannung bei $s=30\,\mathrm{mm}$.
  \begin{ClozeArea}[34]{9}
    Die beiden Messpunkte bestimmen die Empfindlichkeit und den Achsenabschnitt:
    \[
      \begin{aligned}
        m&=\frac{9\,\mathrm{V}-1\,\mathrm{V}}
                 {50\,\mathrm{mm}-10\,\mathrm{mm}}
          =0{,}2\,\frac{\mathrm{V}}{\mathrm{mm}},\\[0.6em]
        b&=1\,\mathrm{V}
          -0{,}2\,\frac{\mathrm{V}}{\mathrm{mm}}\cdot10\,\mathrm{mm}
          =-1\,\mathrm{V},\\[0.6em]
        U(s)&=0{,}2\,\frac{\mathrm{V}}{\mathrm{mm}}\,s-1\,\mathrm{V},
        \qquad U(30\,\mathrm{mm})=5\,\mathrm{V}.
      \end{aligned}
    \]
    Das Modell ist zunächst nur im kalibrierten Messbereich
    $10\,\mathrm{mm}\le s\le50\,\mathrm{mm}$ abgesichert.
  \end{ClozeArea}
\end{ExampleBox}
